% ============================================================================
% 2 STATE OF THE ART AND THEMATIC ANALYSIS  —  the literature review.
% Organized BY THEME, not by paper. Each theme ends with the design decision
% it produced. Kept tight: the 8-page limit includes the bibliography.
% ============================================================================
\section{State of the Art and Thematic Analysis}
\label{sec:related}

We read work published between 2018 and 2025, mostly peer-reviewed and from the
main NLP and machine-learning venues (ACL, EMNLP, NAACL, EACL, COLING, ICLR,
NeurIPS), from reviewed journals (\emph{PNAS}, \emph{Computational
Linguistics}, \emph{Frontiers of Computer Science}, \emph{JMIR Medical
Informatics}), or from CHI and WWW for the human-in-the-loop side. Seven sources
are arXiv preprints, kept only where the work is widely cited and from an
identifiable group, and marked as preprints rather than treated as reviewed.
Sources had to speak to one of four questions: how well prompted models extract
and classify, how their output should be shaped, how a reviewer should be
pointed at the doubtful cases, and how any of it can be checked without gold
data. Eight themes follow.

\paragraph{From supervised pipelines to prompted generalists.}
\citet{xu2024generative} survey generative IE. Generative formulations can
accommodate dynamically specified labels through natural-language descriptions,
avoiding a fixed classifier output layer. Specialized instruction-tuned systems
such as InstructUIE \citep{wang2023instructuie} and targeted-distillation systems
such as UniversalNER \citep{zhou2024universalner} demonstrate that
task-specialized smaller models can outperform general-purpose prompted LLMs on
IE benchmarks. GLiNER \citep{zaratiana2024gliner} makes a related point from the
encoder side: a bidirectional encoder matches LLM prompting at a fraction of
the inference cost.
\textbf{In our project,} fine-tuning was out of scope: we had access to one
hosted model and no permission to train on company data. So we compare only
methods that need no training step.

\paragraph{How far prompting alone gets on extraction.}
\citet{han2023solved} ask whether IE is solved by ChatGPT and answer no:
performance holds on tasks rewarding reasoning and drops on sequence labeling,
where boundaries count. \citet{xie2023empirical} show that decomposing NER into
simpler label-specific questions can substantially improve zero-shot
performance. \citet{sivarajkumar2024empirical} compare six prompt types across five
clinical tasks and find no overall winner, so prompt design must be tested per
task. GPT-NER \citep{wang2023gptner} adds a self-verification step asking
whether an extracted span really belongs to the target type.
\textbf{In our project,} two of these claims are testable, so we test both
(H1, H3).

\paragraph{Document-level scientific extraction.}
SciREX was introduced specifically because sentence- and paragraph-level IE
misses information distributed across scientific documents
\citep{jain2020scirex}. It annotates mentions of Method, Task, Metric and
Dataset as well as coreference, salience and document-level relations. Our
study uses its mention spans as externally supplied reference annotations; it
does not attempt the original paper's coreference or $n$-ary relation tasks.
This distinction prevents a misleading direct comparison between our
pre-annotation score and the SciREX end-to-end relation baseline.

\paragraph{Classification without training data, and label-set size.}
\citet{yin2019benchmarking} turn classification into natural language inference
with each label as a hypothesis, which works for any label set at inference
time but reacts strongly to wording. The embedding alternative encodes text and
label descriptions and selects the nearest representation
\citep{reimers2019sentence}; it is computationally inexpensive, local, and
requires no LLM call, and SetFit \citep{tunstall2022efficient}
shows how much it gains from eight labeled examples per class. MTEB
\citep{muennighoff2023mteb} is the usual way to choose an encoder. The harder
problem is scale: \citet{bhambhoria2023simple} show that flat zero-shot
classification degrades on layered label sets and that splitting the
decision into entailment steps recovers much of the loss, TELEClass
\citep{zhang2025teleclass} adds class-indicative terms to the label set,
and HierPrompt \citep{zhang2025hierprompt} builds hierarchy-aware prototypes.
All three assume that putting hundreds of labels in one prompt does not work.
\textbf{In our project,} the label sets have 4, 10 and several hundred values,
so we can watch this drop happen inside one dataset (H2).

\paragraph{Output format is not a free choice.}
\citet{tam2024speak} show format restrictions measurably hurt reasoning, and
tighter constraints hurt more, though structured formats \emph{help}
classification, where the task is choosing rather than reasoning. Span
extraction additionally needs offsets consistent with the emitted text, which
no JSON schema checks.
\textbf{In our project,} each extraction method exists twice, structured and
free-form, with accuracy and structural validity measured separately.

\paragraph{Whether an agentic step helps.}
\citet{huang2024selfcorrect} study intrinsic self-correction and find
performance usually does not improve and can get worse, because the model has
no independent signal about which parts were wrong. Work on LLMs as judges
documents position, verbosity and self-preference biases pointing the same way
\citep{zheng2023judging}.
\textbf{In our project,} we treated ``agentic helps'' as a hypothesis to test
rather than assume, and built two agentic variants for it (H3).

\paragraph{LLMs as annotators, and why validation is not optional.}
\citet{gilardi2023chatgpt} report ChatGPT beating crowd workers on several
annotation tasks at a fraction of the cost. The careful work came soon after.
\citet{ziems2024transform} test 13 models on 25 benchmarks and find fair
agreement with human coders but no win over fine-tuned models,
\citet{reiss2023testing} shows non-determinism alone can push consistency below
accepted reliability thresholds, and \citet{pangakis2023validation} replicate
27 tasks across 11 datasets and conclude every task needs its own check against
human labels. \citet{baumann2025hacking} quantify the downstream risk: across
13 million annotations feeding 1.4 million regressions, ordinary choices of
model and prompt wording sufficed to flip statistical conclusions. See
\citet{tan2024annotation} for a survey.
\textbf{In our project,} this theme is the reason the project exists. It also
fixes one rule for the tool: every human decision gets stored, because the
record of disagreement is the evidence.

\paragraph{Prioritizing uncertain cases for review.}
A reviewer facing forty plausible suggestions stops reading them properly, so
the tool must rank them. \citet{geng2024survey} survey confidence estimation
and split it into white-box methods reading internal signals such as token
probabilities and black-box methods treating the model as opaque;
\citet{kadavath2022know} show large models carry usable self-knowledge in their
probabilities. The black-box route is sampling: \citet{wang2023selfconsistency}
aggregate several sampled generations and \citet{kuhn2023semantic} refine this
further by grouping samples that mean the same thing before counting them.
\citet{xiong2024express} compare the families and report two findings we relied
on: models are overconfident when stating confidence in words, and white-box
methods beat black-box ones only narrowly. For annotation specifically,
\citet{wang2024collaborative} study how to make human verification effective
rather than mechanical.
\textbf{In our project,} the annotation tool supports multiple confidence
signals for prioritizing review; a systematic comparison of their calibration,
review utility, and additional cost remains future work.

\paragraph{Keeping the human in the loop.}
INCEpTION \citep{klie2018inception} already paired automatic suggestions with
human curation before LLMs; MEGAnno+ \citep{kim2024meganno} splits the work as
we do, with LLM agents proposing and people verifying. For annotation quality
the standard vocabulary is \citet{artstein2008intercoder}.
\textbf{In our project,} we took that two-step design, adapted it to span-level
extraction with character offsets, and logged every accept, relabel and
delete.

\paragraph{The integration considered here.}
Existing work studies individual components of this process, but we did not
find a single task-separated workflow combining reproducible LLM inference,
independently scored extraction and classification, reviewer correction,
provenance, and reusable gold-data export for the target setting considered
here. This project implements and evaluates that specific integration.
